\begin{pastexam}[label={px:nccu-113-dsa-09}]{政大資科所 113 學年度 第 9 題}
（單選；不必說明理由）

(4\%) Assume that all edges have distinct weights. Which of the following statements about minimum
spanning tree is wrong?
\begin{enumerate}[label=\Alph*.]
\item The minimum spanning tree cannot include the edge with the largest weight.
\item The minimum spanning tree must include the edge with the smallest weight.
\item The minimum spanning tree must include the edge with the second smallest weight.
\item The minimum spanning tree is unique.
\end{enumerate}
\end{pastexam}

\begin{pastsoln}
答案是 A。 題目問哪一個是\textbf{錯的}。 圖是連通的簡單圖，權重兩兩不同。

\begin{itemize}
\item[A.] 錯。 最重的邊如果是\textbf{橋}（拿掉就不連通），每棵生成樹都要用它。 反例：路徑 $a$--$b$--$c$，
$w(a,b) = 1$、$w(b,c) = 2$，唯一的生成樹就是整張圖，含最重的邊。 只有最重的邊落在某個環上時，才能說它不在最小
生成樹裡（環性質，資結書定理~24.3）。
\item[B.] 對。 設最輕的邊是 $e = (u, v)$。 取切割 $(\{u\}, V - \{u\})$，$e$ 跨過它，而且是全圖最輕的邊，當然是
跨過這個切割最輕的邊。 權重相異時，由切割性質（資結書推論~24.2），$e$ 在唯一的最小生成樹裡。
\item[C.] 對。 Kruskal 演算法（資結書第~24.2~節）依權重由小到大考慮邊。 第一條一定加入；第二條只有在和第一條形成
環時才會被丟掉，但兩條邊在簡單圖裡不會形成環，所以第二條也一定加入。 Kruskal 得到的是最小生成樹，而且權重
相異時最小生成樹只有一棵，所以它一定含第二輕的邊。
\item[D.] 對。 權重兩兩相異時最小生成樹唯一（資結書定理~24.4）。
\end{itemize}
\end{pastsoln}
